\documentclass[10pt,a4paper]{amsart}
\usepackage[margin=19mm]{geometry}
\usepackage{amsmath,amssymb,mathtools}
\usepackage[scr=rsfs,cal=euler]{mathalfa}
\usepackage{xfrac}
\usepackage[unicode,hidelinks,bookmarks=false]{hyperref}
\newcommand{\ie}{i.e.\ }
\newcommand{\cf}{cf.\ }
\newcommand{\resp}{respectively\ }
\newcommand{\Subsection}{\subsection}
\providecommand{\nobreakdash}{\nobreak\hbox{-}}
\setlength{\emergencystretch}{2em}
\multlinegap0pt
\usepackage{amssymb}
\usepackage{xcolor}
\newtheorem{lemma}{Lemma}
\def\PG{\mathrm{PG}}
\def\cS{\mathcal S}
\title{Segre Varieties and Desarguesian Spreads}
\author{Antonio Cossidente, Giuseppe Marino, Francesco Pavese, Paolo Santonastaso, John Sheekey}
\date{Source: arXiv:2605.22194v1 (CC BY 4.0)}
\begin{document}
\maketitle

\noindent\emph{Source and licence.} Segre Varieties and Desarguesian Spreads. Version arXiv:2605.22194v1 (CC BY 4.0), distributed by arXiv under the Creative Commons Attribution 4.0 International licence, http://creativecommons.org/licenses/by/4.0/, as recorded in the arXiv licence metadata for this exact version and shown on the abstract page https://arxiv.org/abs/2605.22194v1. Full licence notice: \texttt{licenses/arxiv-2605.22194-CC-BY-4.0.txt}.\par\smallskip

\noindent\textit{Editorial scope.} Counted unit: the article's lemma lemma:retta1 (source0.tex lines 494--498). The article's complete original proof of it is delivered with it (source0.tex lines 499--549, one \texttt{proof} environment of 51 lines). No mathematics is written, repaired or re-derived: the statement and the proof are the article's own lines, in the article's own notation, and no retained line is substituted or re-flowed.  No further source-local context is needed: every cross-reference of every retained line resolves inside the two delivered spans.  Nothing was omitted from any retained span, so the package carries no omission record.  No retained line contains a citation, so the wrapper carries no bibliography and no bibliography entry.  No retained line contains a figure environment, an image-inclusion command or drawing code, so no image is delivered and the package declares no asset dependency.  Editorial formatting only, in addition: an AMS \texttt{amsart} preamble that restates the article's own declarations and macros used by the retained lines, namely the environments \texttt{lemma} on separate counters, together with the standard packages those lines need.  The retained environments print in this document as Lemma 1 in order of appearance, whereas the article prints the same items with the numbering of its own full document.  The complete native source of the article as distributed in the arXiv e-print for this exact version is bundled unchanged as \texttt{sources/201119/source0.tex}.
\begin{lemma} \label{lemma:retta1}
Let $\cS$ be a subset of points of $\PG(k-1,q^h)$ such that $\cS$ contains $k+1$ points in general position, and such that for any 
$k+1$ points of $\cS$ in general position, the unique $q$-order 
subgeometry they determine is entirely contained in $\cS$. If a line $\ell$ has at least two points in common with $\cS$, then $|\ell \cap \cS| \ge q+1$. Also, if $\ell_1$ and $\ell_2$ are two intersecting lines such that $|\cS \cap \ell_i| \ge 2$, $i = 1,2$, then $\ell_1 \cap \ell_2 \in \cS$. 
\end{lemma}

\begin{proof}
First, we prove that if a line $\ell$ has at least two points in common with 
$\mathcal S$, then $|\ell \cap \mathcal S| \ge q+1$. 
Let $P_1$ and $P_2$ be distinct points of $\mathcal S$, and let $\ell$ be the line 
of $\PG(k-1,q^h)$ through them. Observe that $\cS$ contains a $q$-order subgeometry of $\PG(k-1, q^h)$. Hence we can find  $P_3,\dots,P_{k+1}$ further $k-1$ points 
of $\mathcal S$ such that no $k$ points of $\{P_1,\dots,P_{k+1}\}$ lie in a 
hyperplane. Then the $q$-order subgeometry determined by $P_1,\dots,P_{k+1}$ 
contains $q+1$ points of $\ell$ and, by hypothesis, is entirely contained in 
$\mathcal S$. Hence, if a line $\ell$ has at least two points in common with 
$\mathcal S$, then $|\ell \cap \mathcal S| \ge q+1$, proving the first part of 
the assertion.

\noindent To prove the second part, let $\ell_1$ and $\ell_2$ be two intersecting lines such 
that $|\mathcal S \cap \ell_i| \ge 2$ for $i=1,2$. Let $P_1,P_2 \in \ell_1 \cap 
\mathcal S$ and let $Q_1,Q_2 \in \ell_2 \cap \mathcal S$. If one of the points 
$P_1,P_2$ lies on $\ell_2$, or one of the points $Q_1,Q_2$ lies on $\ell_1$, then 
$\ell_1 \cap \ell_2 \in \mathcal S$ and the assertion follows immediately. 
Therefore, we may assume that this is not the case.

\noindent Denote by $\theta_q$ the $q$-order subgeometry determined by $k+1$ points of 
$\mathcal S$ in general position; in particular, $\theta_q \subseteq \mathcal S$. 
Let $\Lambda$ be a $(k-4)$-dimensional subspace of $\PG(k-1,q^h)$ disjoint from 
$\langle \ell_1,\ell_2 \rangle_{q^h} \simeq \PG(2,q^h)$ such that 
$\Lambda \cap \theta_q \simeq \PG(k-4,q)$. Let $T_1,\dots,T_{k-2}$ be $k-2$ points 
of $\Lambda \cap \theta_q$ such that no $k-3$ of them lie in a $(k-5)$-dimensional subspace of 
$\Lambda$. Since $T_{k-2}$ and $Q_2$ belong to $\mathcal S$, by the first part of the proof 
the line $\langle T_{k-2},Q_2 \rangle_{q^h}$ contains at least $q+1$ points of 
$\mathcal S$. Hence, we may choose a point 
$T' \in \langle T_{k-2},Q_2 \rangle_{q^h} \cap \mathcal S \setminus 
\{T_{k-2},Q_2\}$. Consider the $k+1$ points
\[
T_1,\dots,T_{k-3},T',P_1,P_2,Q_1,
\]
which are in general position. Let $\theta_q'$ be the $q$-order subgeometry 
determined by these points. By hypothesis, $\theta_q' \subseteq \mathcal S$. Moreover, $\Lambda \cap \theta_q' \simeq \PG(k-4,q)$ and 
$\langle T',P_1,P_2,Q_1 \rangle_{q^h} \cap \theta_q' \simeq \PG(3,q)$; hence,
\[
\Lambda \cap \langle T',P_1,P_2,Q_1 \rangle_{q^h} = T_{k-2} \in \theta_q'.
\]
Similarly, 
$\langle T_{k-2},T' \rangle_{q^h} \cap \theta_q' \simeq \PG(1,q)$ and 
$\langle P_1,P_2,Q_1 \rangle_{q^h} \cap \theta_q' \simeq \PG(2,q)$, which implies
\[
\langle T_{k-2},T' \rangle_{q^h} \cap 
\langle P_1,P_2,Q_1 \rangle_{q^h} = Q_2 \in \theta_q'.
\]

\noindent Therefore, $\ell_1$ and $\ell_2$ are two extended lines of the subgeometry 
$\theta_q'$, and it follows that \[\ell_1 \cap \ell_2 \in \theta_q' \subseteq 
\mathcal S,\] completing the proof.
\end{proof}

\end{document}
